% Epigraph character set (as given in your updated epigraph):
% Athena, Quill, Kestrel, Sable, Raven, Aegis, Aletheia, Irin, Axiom, The Lemma Guy, Cyrene

\PuzzleChapterIntro{The Nine Lenses}{Visual Game \textbullet\ Pairing Invariant \textbullet\ Unique Name}{This is a minimax grid game. You must determine the optimal score one player can force against a perfect opponent by identifying a geometric invariant that limits the maximum sum.}
\Lore{

A scryglass table sits in the Archive Casino: five panes by five, each pane a dark socket in brass.

Two keepers take turns filling the sockets: one can write only \textit{Lumen} and the other can stamp only \textit{Null}.

When the last socket is filled, the glass does not read the whole surface.
It reads only its nine viewing lenses (the $3\times3$ windows), then speaks the brightest total.

A brass ring translates totals into names.

\textit{Speak the name the table speaks.}

}

\Problem

On a $5\times5$ grid, two players alternately fill empty cells.
The first player always writes $1$, the second player always writes $0$.
Exactly one cell is filled per turn, until all $25$ cells are filled.

There are exactly nine $3\times3$ \emph{lenses} (the $3\times3$ subgrids with consecutive rows and columns).
For each lens, compute the sum of the nine entries in it.
Let $A$ be the \emph{maximum} of these nine sums.

\medskip

The ring of names is:

\[
\begin{array}{c|c}
A & \text{Name on the Ring}\\\hline
0 & \text{Athena}\\
1 & \text{Quill}\\
2 & \text{Kestrel}\\
3 & \text{Sable}\\
4 & \text{Raven}\\
5 & \text{Aegis}\\
6 & \text{Aletheia}\\
7 & \text{Irin}\\
8 & \text{Axiom}\\
9 & \text{Cyrene}
\end{array}
\]

The first player wants to make $A$ as large as possible; the second player wants to make $A$ as small as possible.

\VaultDirective{What value of $A$ can the first player guarantee, and what \emph{name} on the ring corresponds to it?}

\SolutionPage
\footnotesize

We show the second player can keep $A\le 6$, and the first player can force $A\ge 6$.
Hence $A=6$ and the accepted name is \textbf{Aletheia}.

\medskip
\textbf{Step 1: The second player can enforce $A\le 6$ by pairing.}

Pair each cell in row $1$ with the cell directly below it in row $2$ (same column),
and pair each cell in row $3$ with the cell directly below it in row $4$.
(Row $5$ is unpaired.)
So there are $10$ disjoint pairs.

\smallskip
\noindent\emph{Strategy.}
Whenever the first player writes $1$ in a cell that belongs to one of these pairs,
the second player immediately writes $0$ in the other cell of the same pair (if it is still empty).
Otherwise the second player plays anywhere.

Then each paired two-cell block contains at most one $1$, hence at least one $0$.

Now take any $3\times3$ lens. It uses three consecutive rows, so it contains either both rows $1$ and $2$,
or both rows $3$ and $4$.
In that two-row block, the lens contains three complete pairs (one in each of its three columns),
so the lens contains at least three $0$'s.
Therefore every lens has sum at most $9-3=6$, so $A\le 6$.

\medskip
\textbf{Step 2: The first player can force $A\ge 6$.}

Label rows $a,b,c,d,e$ (top to bottom) and columns $1,2,3,4,5$ (left to right).

\smallskip
\emph{Move 1.} First plays $1$ at the center cell $c3$.

Let the second player's first $0$ be at some cell $q\ne c3$.
Consider the four ``bands'': top three rows $\{a,b,c\}$, bottom three rows $\{c,d,e\}$,
left three columns $\{1,2,3\}$, and right three columns $\{3,4,5\}$.
Their intersection is exactly $\{c3\}$, so $q$ lies outside at least one band.
First chooses such a band and plays the center-adjacent cell inside it.
By rotation/reflection symmetry, we may assume the chosen band is the bottom three rows $\{c,d,e\}$,
so First plays $1$ at $d3$ on the second move by First, and the first $0$ lies outside rows $\{c,d,e\}$.

Now consider the second player's second move.

\medskip
\textbf{Case 1:} Second does \emph{not} play at $e3$.

Then First plays $1$ at $e3$, so First owns the vertical triple $c3,d3,e3$.
Let $S$ be the $12$ cells in rows $\{c,d,e\}$ but not in column $3$.
At this moment, the first $0$ is outside rows $\{c,d,e\}$ and the second $0$ is not at $e3$,
so at most one $0$ lies in $S$; hence at least $11$ cells of $S$ are still empty.

From now on, whenever any cell of $S$ is empty, First plays in an empty cell of $S$.
Even if Second also plays in $S$ whenever possible, First claims at least $\lfloor 11/2\rfloor=5$ cells of $S$.
Split $S$ into $S_L=\{c,d,e\}\times\{1,2\}$ and $S_R=\{c,d,e\}\times\{4,5\}$, each of size $6$.
By pigeonhole, one of $S_L,S_R$ contains at least $3$ of First's ones.
Together with the triple in column $3$, this gives a $3\times3$ lens with at least $6$ ones.
So $A\ge 6$ in Case 1.

\medskip
\textbf{Case 2:} Second \emph{does} play at $e3$ (placing a $0$ there).

Then First plays $1$ at $c2$.
Define the two $y$-cells to be $d2$ and $e2$, and define the six $z$-cells to be
\[
c1,d1,e1,\quad c4,d4,e4.
\]
All eight of these cells lie in rows $\{c,d,e\}$ and are distinct from $c2,c3,d3,e3$,
so all eight are empty at this moment.

First now plays by priorities:
\begin{itemize}
\item If Second ever plays in one $y$-cell, then on the next move First plays in the other $y$-cell.
\item Otherwise, whenever any $z$-cell is empty, First plays in an empty $z$-cell.
\item If no $z$-cell is empty and both $y$-cells are still empty, then First plays in a $y$-cell.
\end{itemize}
This guarantees that First gets at least one $y$-cell.
Also, among the six $z$-cells, Second can claim at most three, so First claims at least three $z$-cells.
Since three $z$-cells lie in column $1$ and three in column $4$, having at least three $z$'s forces
First to have at least two $z$'s in the same one of these columns.

If First has at least two $z$'s in column $1$, then the bottom-left lens $\{c,d,e\}\times\{1,2,3\}$ contains
$c2,c3,d3$ plus those two column-$1$ ones plus one $y$-cell in column $2$, totaling at least $6$ ones.
If instead First has at least two $z$'s in column $4$, then the lens $\{c,d,e\}\times\{2,3,4\}$ has the same conclusion.
So $A\ge 6$ in Case 2 as well.

\medskip
Combining both cases, First can force $A\ge 6$.
Together with Step 1, this gives $A=6$.

On the ring, $A=6$ corresponds to \textbf{Aletheia}.

\FinalAnswer{Aletheia}

\NextPage
